
The best-performing method for spurious correlation robustness does not change the backbone at all.

Deep Feature Reweighting (DFR) \citep{kirichenko2022last} achieves WGA = 0.91 on Waterbirds --- outperforming GroupDRO (0.88), SAM (0.74), and ERM (0.72) --- yet its ResNet-50 backbone is numerically identical to the ERM backbone it reuses (layer4 cosine similarity = 1.000000 $\pm$ 1e-14 across all 3 matched seed pairs). DFR retrains only the classification head on group-balanced held-out data while freezing every backbone weight. GroupDRO \citep{sagawa2019distributionally}, which achieves WGA = 0.88, modifies backbone weights substantially (backbone/head L2 ratio = 6.47--7.03). Both methods improve substantially over ERM (WGA = 0.72), yet through fundamentally different interventions.

This structural situation --- one method leaving the backbone identical to ERM while achieving the highest WGA; another achieving near-best WGA by substantially modifying backbone weights --- reveals a gap in mechanistic understanding of robustification methods. The field has focused on \textit{what} these methods achieve (WGA improvement) rather than \textit{how} --- specifically, whether improvement comes from changing the backbone's spurious feature encoding, recalibrating the head to exploit existing features differently, or both.

The key enabling insight is that backbone-level spurious feature encoding and head-level decision making can be measured separately. Freezing the backbone and training a linear probe on frozen layer4 features to classify the spurious attribute (land vs. water background) produces a direct, falsifiable measurement of backbone spurious encoding --- independent of head weights. Methods that reduce backbone spurious encoding should produce lower probe accuracy; methods achieving WGA improvement through head recalibration alone should not affect probe accuracy.

We formalize and test this diagnostic framework with pre-registered statistical tests on 12 publicly available ResNet-50 checkpoints (izmailovpavel/spurious\_feature\_learning \citep{izmailov2022feature}, 3 seeds $\times$ 4 methods). Our analysis confirms a three-step mechanistic chain for GroupDRO: minority group upweighting (H-M1) $\rightarrow$ group-balanced gradient propagation through all backbone layers (H-M2) $\rightarrow$ significantly reduced background linear decodability in layer4 features (H-M3: p = 0.0039, Cohen's d = 6.48). Simultaneously, H-P0 establishes that DFR backbone features are mathematically identical to ERM backbone features --- the highest WGA method leaves the backbone unchanged.

\textbf{Contributions.} This diagnostic study makes the following contributions:

\begin{enumerate}
\item \textbf{Empirical backbone-vs-head typology.} GroupDRO and DFR represent two structurally distinct robustification pathways: backbone-modification (GroupDRO, measurably reducing layer4 spurious encoding, p = 0.0039, d = 6.48) versus head-recalibration (DFR, backbone cosine similarity = 1.000000 with ERM). This typology is established with pre-registered statistical tests against 12 publicly available checkpoints.
\end{enumerate}

\begin{enumerate}
\item \textbf{Quantified backbone divergence between methods.} GroupDRO modifies backbone weights 6.47--7.$03\times$ more than head weights (backbone/head L2 ratio), while DFR-ERM weight differences are exactly 0.000 for all backbone layers across all seeds. This is the first per-seed, per-method measurement of this distinction for the izmailovpavel checkpoints, filling the measurement gap in Izmailov et al. [2022].
\end{enumerate}

\begin{enumerate}
\item \textbf{Mechanistic chain for GroupDRO robustification.} All three steps of the proposed mechanism are verified with pre-registered gates: minority group upweighting (minority fraction = 5.01\%), gradient propagation to backbone (backbone/head ratio = 6.47--7.03), and reduced spurious decodability (ERM 0.9838 $\rightarrow$ GroupDRO 0.9530; p = 0.0039, d = 6.48).
\end{enumerate}

\begin{enumerate}
\item \textbf{WGA-spurious encoding correlation (exploratory).} Background probe accuracy correlates negatively with WGA across 9 checkpoints (r = $-0.504$, SUGGESTIVE; ERM+GroupDRO subset r = $-0.755$, p = 0.041, CONFIRMED), suggesting backbone spurious encoding level is a meaningful correlate of WGA performance.
\end{enumerate}

The paper proceeds as follows: Section 2 reviews related work. Section 3 describes the diagnostic methodology. Section 4 details experimental design. Section 5 presents results. Section 6 discusses implications and limitations. Section 7 concludes.

